Se pide lo siguiente en relación con el precio sombra de la restricción $i$-ésima de un problema de programación lineal:

\begin{enumerate}[label=\alph*)]
\item \textbf{Indicar cuál es su expresión matemática.}

El precio sombra se define como $\pi^B=c^BB^{-1}$ y la componente $i$-ésima es $\pi^B_i$.

\item \textbf{Enunciar cuál es su interpretación y qué información aporta.}

El precio sombra de la restricción $i$ es igual al incremento de la función objetivo cuando la disponibilidad del recurso correspondiente a la restricción $i$ aumenta una unidad, es decir:

\[
\pi^B_i = \Delta z|_{\Delta b_i=1}
\]

\item \textbf{Demostrar razonadamente qué información aporta.}
La expresión del valor de la función objetivo de un problema es:

\[
z = cx=c^Bu^B=c^BB^{-1}b=\pi^Bb
\]

Explicitando la componente $i$-ésima de $\pi^B$:

\[
z = \begin{pmatrix}
   \pi^B_1 &  \pi^B_2  & \cdots & \pi^B_i & \cdots & \pi^B_m
   \end{pmatrix}
   \begin{pmatrix}
   b_1 \\
   b_1 \\
   \cdots \\
   b_i \\
   \cdots \\
   b_m \\
   \end{pmatrix}
   = \pi^B_1b_1 + \pi^B_2b_2 + \cdots + \pi^B_ib_i + \cdots +\pi^B_mb_m
\]

Si aumenta una unidad $b_i$, la expresión anterior sería:

\begin{align}
z\big|_{\Delta b_i = 1} &= \pi^B_1 b_1 + \pi^B_2 b_2 + \cdots + \pi^B_i (b_i + \mathbf{1}) + \cdots + \pi^B_m b_m  \notag \\ 
&= \pi^B_1 b_1 + \pi^B_2 b_2 + \cdots + \pi^B_i b_i + \cdots + \pi^B_m b_m + \mathbf{\pi^B_i} \notag \\
&= z + \pi^B_i \notag
\end{align}

De modo que se aprecia que la función objetivo ha aumentado en un valor igual a $\pi^B_i$



\end{enumerate}